% Task 2: the text of the task. Tasks.tex puts all tasks into one PDF; the code shown here is in Task2/.
\settitle{Task 2}{The discount code}

\begin{story}
An online shop gave a customer the discount code \texttt{PROMO50}. The code may be used \textbf{at most 3 times}:
the shop checks that in the code, with a counter. A week later the accountants saw that the code had been used
several hundred times, although the limit was 3. How is this possible?
\end{story}

\section{The situation}

Every click on \emph{Use code} starts a new thread that calls \texttt{redeem()}. The function checks the counter
\texttt{used}, verifies the order (slow: stock, price, address) and counts the use.

\begin{center}
\begin{tikzpicture}
  % shoppers with phones
  \foreach \i/\y in {0/1.5, 1/0.5, 2/-0.5, 3/-1.5} {
    \node[text=sIndigo, font=\LARGE] (ph\i) at (0,\y) {\faMobile*};
    \node[draw=sIndigo, fill=sIndigo!12, rounded corners=2pt, font=\scriptsize\ttfamily, text=nInk, anchor=west,
      inner sep=2pt] (btn\i) at (0.4,\y) {Use PROMO50};
    \draw[-{Stealth}, thick, sIndigo] (btn\i.east) -- (3.35,\y*0.35);
  }
  \node[font=\scriptsize, text=nInk] at (1.1,-2.25) {100 clicks at once};
  % the shop
  \node[hw, minimum width=4.6cm, minimum height=2.6cm] (shop) at (5.7,0) {};
  \node[text=nInk, font=\Large, anchor=north west] at ($(shop.north west)+(0.2,-0.15)$) {\faShoppingCart};
  \node[font=\bfseries, text=nInk, anchor=north west] at ($(shop.north west)+(0.95,-0.22)$) {online shop};
  \node[font=\scriptsize\ttfamily, text=nInk, anchor=west, align=left] at ($(shop.west)+(0.25,-0.35)$)
    {1. used < 3 ?\\2. verify\_order() \textrm{(100\,ms)}\\3. used++};
  % what the shop gave out
  \foreach \k in {0,...,11} {\pgfmathsetmacro{\xx}{9.0+mod(\k,4)*0.62} \pgfmathsetmacro{\yy}{0.9-floor(\k/4)*0.62}
    \node[text=sPink, font=\large] at (\xx,\yy) {\faTag};}
  \node[font=\scriptsize, text=nInk] at (9.95,-1.0) {$-50\%$ discounts given};
  \draw[-{Stealth}, thick, nInk] (shop.east) -- (8.6,0);
  \node[font=\small\bfseries, text=AccentOrange, align=center] at (9.95,-1.75) {limit: 3\\given: hundreds};
\end{tikzpicture}
\end{center}

\begin{note}{Ready code (in \texttt{coupon.h}, we do not look inside)}
\textbf{Functions}\\[2pt]
\begin{tabularx}{\linewidth}{@{}p{3.6cm}X@{}}
\texttt{verify\_order()} & simulates the shop checking the order (stock, price, address): 100\,ms \\
\end{tabularx}
\par\smallskip\textbf{Constants}\\[2pt]
\begin{tabularx}{\linewidth}{@{}p{3.6cm}X@{}}
\texttt{MAX\_USES} & 3: the code may be used at most 3 times \\
\end{tabularx}
\end{note}

This bug is not only a wrong number in a report. Here it costs the shop money, and the same mistake in a
check of a limit (a number of tries, a budget, a quota) is a \textbf{security hole}. Programmers call it
\emph{time of check to time of use} (TOCTOU): the program checks a rule at one moment and uses the result
a moment later, when the rule may be already broken.

\codesection

The template is in \ghfolder{Task2}. This is \texttt{coupon.cpp}:

\codefile[firstline=6]{coupon.cpp}

\runline

\section{Your tasks}

\begin{questions}
  \item Does the limit of 3 uses really work? Suppose the code has already been used 2 times, so one use is left.
        Two customers click \emph{Use code} at the same moment. Describe, step by step, how both of them can get
        the discount.
  \item Set \texttt{CLICKS} to 2, 10, 100 and 1000. Run each one 3 times and write the results in a table.
        How many times was the code used? What do you notice?
  \item Fix \texttt{redeem()} so that the code works at most 3 times, for any number of clicks.
  \item \texttt{verify\_order()} is slow, and while one click verifies its order under the mutex, all the others
        wait. It is possible to move \texttt{verify\_order()} \textbf{outside} the mutex and still allow at most
        3 uses. Think how to do it (Task 1, question 4 is similar) and write it.
\end{questions}
